En un estanque se modela la cantidad de un organismo mediante su biomasa $Q(t)$, medida en gramos. La producción por reproducción y la pérdida por mortalidad son proporcionales a $Q$, con tasas $1.8$ y $1.4\,\mathrm{año}^{-1}$, respectivamente. Además se añaden $50\,\mathrm{g/año}$. Inicialmente hay $50\,\mathrm g$.

\begin{enumerate}[label=(\alph*),ref=\alph*,leftmargin=*]
\item Formule la ecuación y el dato inicial.
\item Obtenga $Q(t)$ y calcule $Q(10)$.
\end{enumerate}
